% !TEX root = tarasov-qasp-zh.tex
% 尾页：tikz 框架照抄 00-back.tex；原文件书名与社名所用 \textls[150]（字距拉伸）
% 在 xelatex 下不可用，按全书体例改为 \textbf（书名、社名保留英文）；说明段用前言定译。
\thispagestyle{empty}

\begin{tikzpicture}[remember picture,overlay]
  \fill[thick,yellow!10]
        (current page.north west) rectangle (current page.south east);
\end{tikzpicture}
\begin{tikzpicture}[remember picture,overlay] 

\node [align=justify,scale=1.25,yshift=2cm,xshift=0cm,font=\sffamily,Maroon,text width=\textwidth]	 at (current page.center) {本书是为准备高校物理入学考试的学生编写的辅助读物，采用作者（\textbf{教师：}）与好问的读者（\textbf{学生：}）对话的形式。这对于分析考生在入学考试中的常见错误、复习同一问题的多种解法、讨论物理理论中的疑难问题，都特别方便。
};

\end{tikzpicture}

\begin{tikzpicture}[remember picture,overlay] 
\node [scale=3.,align=left,font=\sffamily,yshift=2.5cm,xshift=-0.09cm,SteelBlue] at (current page.center) {\textbf{QUESTIONS} \\ \textbf{AND ANSWERS} \\ \textbf{IN SCHOOL PHYSICS}};

\end{tikzpicture}

\begin{tikzpicture}[remember picture,overlay] 
\node [yshift=-8.75cm]	 at (current page.center) {\includegraphics[width=0.15\textwidth]{figs/front-back/mir-logo.pdf}};
\node [scale=2.25,yshift=-4.8cm,font=\sffamily,SteelBlue]	 at (current page.center) {\textbf{Mir Publishers Moscow}};
\end{tikzpicture}
